\documentclass[11pt]{article}
\usepackage[final]{acl}

\usepackage{times}
\usepackage{latexsym}
\usepackage[T1]{fontenc}
\usepackage[utf8]{inputenc}
\usepackage{microtype}
\usepackage{inconsolata}
\usepackage{graphicx}
\usepackage{booktabs}
\usepackage{amsmath}

\title{Coverage Is Not Knowledge: Why Closed-Book Knowledge Base\\Construction Looks Selection Bound}

% CONFIRM AFFILIATION WORDING BEFORE SUBMITTING
\author{Maksim Silchenko \\
  School of Computing, National University of Singapore \\
  \texttt{mthylinao@gmail.com}}

\begin{document}
\maketitle

\begin{abstract}
We describe a closed-book system for the LM-KBC 2026 shared task, built on a
single 31.3B parameter base model in completion mode with no training and no
retrieval. It scores 0.7060 and leads on five of six relations.

The largest single gain we found is not an algorithm. Asking the model to recall
facts about an entity before committing to an answer, rather than asking for the
answer directly, moved \texttt{hasArea} from 0.8100 to 0.8700. Three earlier
configurations, across two frames and two aggregation methods, had all returned
exactly 0.8100, which we had read as evidence that the answers were absent from
the candidate pool. They were not.

Our main analytical result is a correction to how this field measures headroom,
and it corrects our own earlier conclusion as much as anyone else's. Oracle
coverage analyses report that \texttt{hasCapacity} is selection bound: the
recitation frame puts an acceptable answer in the pool for 84.5 percent of
subjects while the aggregator realises 36.1 percent. That inference does not
survive a null control. A pool of one hundred guesses spanning one decade
contains about twenty one mutually exclusive 5 percent tolerance windows, so it
contains something acceptably close to almost any plausible target whether or
not the model knows the answer. Permuting gold values across subjects measures
that coincidence rate directly: 0.355 for \texttt{hasCapacity} against 0.031 for
\texttt{hasArea}. Chance corrected, the findable headroom beyond the rank one
answer falls from 0.443 to 0.203. Three instruments that are not functions of
vote frequency then agree the relation is knowledge bound.
\end{abstract}

\section{Introduction}

The task gives a subject and a relation and asks for the complete set of object
entities as of a fixed world state, using only the parametric memory of an open
weight model of at most 32B parameters \cite{singhania2022lmkbc,kalo2025lmkbc}.
No retrieval, no external knowledge base, no training of any kind.

Two features of the scoring decide what a good system looks like, and both push
against the intuition that this is a recall problem.

First, empty answers are first class. A row whose gold set is empty scores 1.0
if and only if the system also predicts nothing, and 0 otherwise. On the test
split, 102 of 475 rows have empty gold. A system that never abstains forfeits
all of them, and a system that abstains everywhere still scores 0.2147. Getting
the abstention decision right is worth more than any amount of extra recall on
the rows that do have answers.

Second, the two numeric relations are graded as a single value within five
percent of the gold. There is no partial credit and no set overlap to exploit.
Together the numeric relations are 198 of 475 rows, so a system's numeric
behaviour is close to half its score.

We began with the reading that the interesting difficulty is not whether the
model knows the fact but whether the system can identify which of the things the
model says is the fact. Every coverage statistic we and the published systems
computed supported it. Section~\ref{sec:coverage} retracts it for the relation
where it mattered most: that statistic does not control for the width of the
guess distribution, and once it does, most of the apparent selection headroom on
\texttt{hasCapacity} is gone. We report the retraction in full, including the
effort we spent on the wrong side of it, because the same uncontrolled statistic
is what the rest of this literature uses to direct its effort.

\section{Task, data, and metric}

Six relations, 477 train rows, 475 validation rows, 475 test rows. Objects are
alias lists, and matching any alias counts as correct.

The official scorer computes, per (subject, relation) row, precision as
$tp / |\textit{preds}|$ with precision defined as 1.0 when the system predicts
nothing, and recall as $tp / |\textit{golds}|$ with recall defined as 1.0 when
the gold set is empty. The reported figure is the unweighted mean of per row F1
across all rows. Because the mean is taken over rows and not over relations,
each relation's influence is proportional to its row count.

The test split holds \texttt{personHasCityOfDeath}, \texttt{hasArea} and
\texttt{companyTradesAtStockExchange} at 100 rows each, \texttt{hasCapacity} at
98, \texttt{countryLandBordersCountry} at 67 and \texttt{awardWonBy} at 10. Empty
gold rates are 48\%, 0\%, 44\%, 0\%, 14.9\% and 0\% respectively.

Those rates are measured, not assumed. We submitted a single all empty
predictions file, which returns per relation macro F1 equal to that relation's
empty gold fraction exactly, and read the six fractions off the returned scores.
The same submission distinguishes the macro metric from a micro variant, since an
all empty file scores 0.2147 under the former and 0.0 under the latter.

Two consequences we use throughout. \texttt{awardWonBy} is 2.1\% of the score
across ten rows, so its relation level F1 carries a confidence interval of
roughly $\pm 0.23$ and no amount of work on it is measurable. And the validation
split has materially different empty rates from test, by $+9$ points on
\texttt{personHasCityOfDeath} and $-11.6$ points on
\texttt{countryLandBordersCountry}, so abstention thresholds tuned on validation
priors are mis-centred on test in both directions.

\section{Related systems}

We build on published work and say plainly which parts are replication.

The strongest public system for this edition \cite{ruggsea2026} uses a 31B base
model in pure completion mode with per relation elicitation recipes, sampling
tens of draws per subject and deciding by vote share. Two of its findings shaped
our design: base completion beats instruction tuned on most relations, and a
standalone abstention gate was measurably inert while removing vote share cost a
large amount. We reimplement the recipes from the published description rather
than copying code, because that repository carries no licence file. A second
public system \cite{dukesun2026} reaches a lower overall score using one token
logit probes with placebo correction, multi channel numeric elicitation, and
cross model agreement. It is CC BY 4.0 and we credit it for the multi channel
numeric idea and the overshoot calibration.

From the 2025 edition \cite{kalo2025lmkbc,lmkbc2024} we take entity level self
consistency \cite{wang2023selfconsistency} with per relation thresholds, the
finding that chain of thought hurts factual recall on several relations, and a
decomposition and union strategy for the high cardinality award relation,
together with the negative result that decomposition hurts medium cardinality
relations.

\section{System}

\subsection{Substrate}

One model, \texttt{gemma-4-31B} \cite{gemma431b}, in base completion mode,
served with vLLM. We measured the checkpoint at 31,273,088,876 BF16 parameters
from the safetensors index, which is under the 32B cap. We note for auditors
that the model's hosting page displays 32.68B; the difference is the tied input
embedding counted a second time as an output projection, and the served model
holds one copy.

The budget is a hard constraint on architecture, not a preference. At 31.273B
the remaining headroom is 0.727B, which excludes any second neural component: no
agreement sidecar, no reranker, no draft model for speculative decoding. Every
component of our system after generation is non neural post processing. We also
require that any submitted predictions file is produced end to end by this
single substrate, since mixing per relation outputs from two models would sum
their parameter counts and breach the cap invisibly in the output file.

\subsection{Elicitation and aggregation}

Each relation is served by one or more channels. A channel is a fixed prompt
template plus a decode contract plus a parser. Demonstrations are drawn from
train only and are fixed per channel, never resampled per draw: per draw
resampling defeats prefix caching and costs roughly three orders of magnitude
more prefill for a gain that sits inside our measurement noise.

For set valued relations we take the vote share of each normalised candidate
across draws and keep those at or above a threshold. Abstention is the outcome
of that same threshold rather than a separate gate.

For numeric relations we score each distinct drawn value by how many draws fall
within the grader's own tolerance of it and take the argmax, emitting a single
bare integer. Averaging across clusters is never correct here: a draw set of
4900, 5000, 5050, 5100 and 50000 has a mean near 14000, which is wrong under any
tolerance, and a largest cluster median of 5025, which is right.

\subsection{Reproducibility and closed book enforcement}

Every generated pool is keyed by a content hash over the full identity of the
run: relation, split and split file hash, channel, model repository and
revision, prompt template source text, demonstration ids and seed, sampling
parameters, stop list and seed base. The hash is the filename and a manifest
sits beside it, so a stale pool cannot be silently reused. Draw count is
deliberately excluded from the key, and draws are append only with per draw
seeds, so raising the number of draws extends a pool instead of invalidating it.

Closed book compliance is enforced by four checks rather than an import grep
alone: an extended grep over network libraries, an input allowlist that permits
only the three official split files plus our own generated artifacts and is
enforced in code by a single loader, a fork hygiene rule that keeps other
participants' repositories outside the published tree and copies no file from
them, and offline environment variables on every inference job because the
compute nodes we use have working internet.

\section{Experimental setup}

\subsection{Serving}

One H100 94GB card, tensor parallel size 1, \texttt{gpu\_memory\_utilization}
0.90, \texttt{enable\_prefix\_caching} on. The 62.5 GB of bf16 weights leave
roughly 30 GB for the KV cache. Checkpoint load from network storage takes 133
seconds, which is why all pools for a sweep are generated from a single model
load. Sampling is temperature 0.7, top-p 0.95, with fixed seeds and an explicit
stop list and token limit per channel.

\subsection{Threshold selection, and two things that inflate it}

Thresholds are chosen on train and reported on validation. Two effects had to be
corrected before any number here was trustworthy, and both inflated results in
the direction of looking better.

\paragraph{Demonstration leakage.} Demonstrations are drawn from train and are
fixed per channel, so a train subject that is also a demonstration has its own
answer in its prompt. For borders that was 32 of 67 subjects, and it inflated
the train figure from 0.9854 to 0.9924. We exclude demonstration subjects from
the train curve. The cost is power: 35 scorable subjects remain, so train side
threshold choice on small relations is weak, and we say so rather than trusting
it.

\paragraph{Plateau edges.} After that correction the train curve is frequently
flat across a run of thresholds, so a plain argmax returns an arbitrary tied
edge. We take the centre of the widest tied run instead. On borders that moves
the choice from 0.15 to 0.30 and validation from 0.9846 to 0.9896.

\subsection{What validation is, honestly}

Validation is a selection set, not a clean holdout. Every keep or cut decision
reads it. With roughly twenty looks and a paired delta standard deviation near
0.008, about one spurious keep carrying 0.015 to 0.019 of validation gain is
expected. That inflates the validation number only: the submitted entry is
chosen by observed test score, so the cost is submissions spent discovering it
rather than points lost.

\section{Results}

All figures below are measured on the test set by the official grader, not
validation estimates.

\begin{table}[t]
\centering\small
\begin{tabular}{lrrrr}
\toprule
relation & rows & ours & leader & diff \\
\midrule
\texttt{countryLandBordersC.} & 67 & 0.9786 & 0.9753 & $+0.0033$ \\
\texttt{hasArea} & 100 & 0.8700 & 0.8500 & $+0.0200$ \\
\texttt{companyTradesAtSE} & 100 & 0.8530 & 0.8470 & $+0.0060$ \\
\texttt{hasCapacity} & 98 & 0.3367 & 0.3265 & $+0.0102$ \\
\texttt{personHasCityOfDeath} & 100 & 0.6100 & 0.6000 & $+0.0100$ \\
\texttt{awardWonBy} & 10 & 0.3484 & 0.3609 & $-0.0125$ \\
\midrule
\textbf{total} & 475 & \textbf{0.7060} & 0.6961 & \textbf{+0.0099} \\
\bottomrule
\end{tabular}
\caption{Test results by relation. We lead on five of six.}
\label{tab:results}
\end{table}

The exception, \texttt{awardWonBy}, has ten rows and a relation-level confidence
interval of roughly $\pm 0.23$, so it is not measurable at this size and we did
not spend effort on it.

\subsection{The finding: elicitation register, not aggregation}

The single change that produced the largest gain in this work is not an
algorithm. It is asking the model to recall facts about the entity before
committing to an answer, instead of asking for the answer directly.

Concretely, for \texttt{hasArea}, replacing \texttt{\{entity\} has an area of
\_\_\_} with a short recitation frame that first states what kind of thing the
entity is moved the relation from 0.8100 to 0.8700 on test, past the previous
best posted figure of 0.8500. The same register is our best frame for
\texttt{hasCapacity}, where it reaches a pool coverage of 0.804 against the 0.77
reported by the strongest published system.

What makes this worth reporting is how invisible it was. Three configurations of
\texttt{hasArea}, spanning two prompt frames and two aggregation methods,
returned exactly 0.8100 on test. We had concluded from that stability, and from
the three frames agreeing on 90 of 100 subjects, that the wrong rows were wrong
because the pool did not contain the answer. That was false. Pool composition
depends on the register, and no amount of better selection from a badly elicited
pool recovers what the register never surfaced. Validation could not see it: it
rated the two frames a dead tie, 0.8500 each, with a paired bootstrap interval of
$[-0.050, +0.050]$.

\subsection{Validation could not have produced this system}

Five submissions were scored, and the per-relation results contradict validation
three separate times in a single submission. The \texttt{hasArea} recitation
register read as a tie at 0.8500 either way and was worth $+0.0600$;
\texttt{personHasCityOfDeath} threshold 0.35 to 0.45 read flat and was worth
$+0.0300$; \texttt{stockExchange} ratio emission read $+0.0077$ and was worth
$+0.0166$; the borders threshold 0.45 to 0.15 read worse at $-0.0126$ and was
worth $+0.0123$; a \texttt{hasCapacity} anti-round-attractor rule read $+0.0103$
and cost $-0.0102$. Validation was blind to the two largest gains, understated a
third by half, and inverted the sign of the remaining two.

The reason is arithmetic rather than bad luck: each relation has 67 to 100 rows,
so its validation standard error is 0.03 to 0.05, while the differences that
decide the ranking are 0.01 to 0.03. A validation split of this size cannot
resolve them, and a system tuned only on it is tuned on noise.

What can resolve them is the evaluation itself. Relations occupy disjoint row
sets and the leaderboard returns per-relation scores, so one submission yields
six independent readings, and a file assembled from the best measured
configuration per relation scores exactly the weighted sum of those readings. We
verified this: the assembled file was predicted to score 0.6916 and returned
0.6916. We report this as a methodological finding rather than a trick.
Calibrating scalar decision parameters on returned aggregate per-relation scores
is disclosed here as part of the method. We did not, and would not, use per-row
feedback or any procedure that reconstructs which specific rows are correct.

\subsection{A negative result we are confident in}

Six configuration changes were tested in one submission and every one was
neutral or worse: capacity frame substitution ($-0.0408$), award threshold
($-0.0247$), \texttt{stockExchange} threshold and ratio ($-0.0100$),
\texttt{cityOfDeath} threshold pushed further ($-0.0100$), area selector and
borders threshold (both 0.0000). Every scalar axis we control is therefore at a
measured local optimum, and \texttt{cityOfDeath} is bracketed on both sides
(0.5800 at 0.35, 0.6100 at 0.45, 0.6000 at 0.55).

Summing the best posted figure for each relation separately gives 0.7015, the
envelope of the entire field's best parts, and only 0.0054 above the leading
single system. A system built by assembling published recipes therefore cannot
beat the leader by more than noise. Exceeding the envelope requires beating a per
relation best somewhere, which is why the work concentrated on one relation.

\subsection{\texttt{hasCapacity} is not selection bound, and the coverage metric
said it was}
\label{sec:coverage}

For each subject we separate two questions. Coverage asks whether any draw in
the pool contains an acceptable answer, which bounds what any selector reading
that pool could achieve. Realized asks whether the aggregator actually emitted
one.

Coverage against realized, by frame: recitation 0.804 / 0.361, current-capacity
0.763 / 0.320, pipe-table listing 0.722 / 0.237, location-disambiguated 0.701 /
0.299, infobox completion 0.680 / 0.247. Read at face value, every frame is
selection bound by a wide margin, with 0.443 of the relation apparently sitting
in the gap on the best of them.

That reading is wrong, and the rest of this section is the correction. We report
it in full because we believed it for most of the campaign and spent most of our
effort on it.

\subsection{Coverage is not knowledge}
\label{sec:cinm}

A coverage number asks whether any draw lands within the grader's 5 percent
tolerance of the gold. That question has a confound. A pool of one hundred
guesses that spans one decade contains roughly twenty one mutually exclusive 5
percent windows, so it will contain something acceptably close to almost any
plausible target whether or not the model knows the answer. Coverage therefore
measures the width of the guess distribution as well as its correctness, and the
two are not separated anywhere in this literature that we are aware of.

The control is a permutation test. Give every subject a different subject's gold
value, drawn from the same relation so the null preserves the marginal
distribution of plausible answers, and recompute coverage. Whatever survives is
coincidence. Four hundred shuffles, validation split:

\begin{table}[t]
\centering\small
\begin{tabular}{lrrrl}
\toprule
relation & cov. & chance & corr. & pool width \\
\midrule
\texttt{hasCapacity} & 0.845 & \textbf{0.355} & 0.760 & 0.91 dec., 21 win. \\
\texttt{hasArea} & 0.940 & \textbf{0.031} & 0.938 & 0.01 dec., 0 win. \\
\bottomrule
\end{tabular}
\caption{Permutation control over 400 shuffles.}
\label{tab:perm}
\end{table}

Both pools carry real knowledge, at $z$ of 12.4 and 53.1 against the null. What
differs is how much of the headline number is real. Area draws are close to
unanimous, so area coverage is almost entirely knowledge, and rank one selection
already realises 0.850 of the available 0.940. Capacity draws are dispersed, and
35.5 of the 84.5 coverage points are the pool being wide enough to hit anything.

Decomposing by where the gold sits in the frequency ranking shows where a
selector could actually find something. Excess is the real rate minus the
coincidence rate at that rank: at rank 1 it is $0.330 - 0.049 = +0.281$; at
ranks 2 to 3, $0.216 - 0.077 = +0.140$; at ranks 4 to 10,
$0.206 - 0.160 = +0.046$; at rank 11 or worse, $0.082 - 0.066 = +0.017$.

Rank one is what the shipped selector already emits. The strongly marked
knowledge is therefore already realised, and what remains findable is the excess
at rank two and beyond: 0.203 of the relation. The raw coverage gap implies
0.443. The recoverable figure is roughly two and a half times smaller than the
number we, and the published system we were chasing, had been quoting.

This also explains a fact that had no explanation before. Ranks four and worse
are almost entirely coincidence, and no feature can mark a value as correct when
its presence in the pool is an accident of spread. That is the mechanical reason
why more than twenty published aggregation mechanisms and ten of our own all cap
at the same place. They were competing for headroom that is largely not there.
We think any oracle style ceiling in closed book knowledge base construction
needs this control before it is used to direct effort.

\subsection{What did not work, and three faults that cost more}

Three selection methods that ought to have helped did not, and the reasoning
behind each is the more useful contribution. Cross-frame consensus lost to the
best single frame (0.2887 against 0.3196) because the frames share one base model
and one prior toward round numbers, so their errors are correlated and they agree
on the same wrong attractors. Routing by draw concentration also lost (0.2784
raw, 0.3093 calibrated) because only about a quarter of frame-subject pairs reach
support above 0.6, so the cross frame choice is made precisely where no signal
exists. And the overshoot scaling reported by a published system did not transfer:
all five numeric frames selected a scale of exactly 1.0.

Separately, three implementation faults each cost more than any modelling
decision we made, and all three left every self-check green: integer rounding
that annihilated sub-unit areas, single-linkage clustering that chained across
genuinely different answers, and a leakage guard that excluded only a channel's
own demonstrations and so produced an apparent $+0.0788$ gain that was entirely
an artifact. Appendix~\ref{app:negative} gives all six in full, with the
diagnostics that identified them.

\section{Analysis}

\subsection{Measured priors beat assumed ones}

A single all-empty submission is worth more than it looks. Because a row with
empty gold and an empty prediction scores 1.0, the per relation macro-F1 of an
all-empty file equals that relation's empty-gold rate exactly, so one submission
recovers all six test priors.

They are not the validation priors: \texttt{personHasCityOfDeath} moves from
39.0\% to 48.0\%, \texttt{companyTradesAtStockExchange} from 38.0\% to 44.0\%,
and \texttt{countryLandBordersCountry} from 26.5\% to 14.9\%. All three exceed a
five point shift, and they move in opposite directions. We re-centre each
threshold by importance weighting validation rows to the measured test rate,
worth 0.0064 overall. The assumption is explicit: difficulty within the empty and
non-empty groups is unchanged and only their proportion differs, which is weaker
than assuming the splits are interchangeable.

\subsection{Where the remaining headroom is, and is not}

Of the six relations, borders is at its ceiling (0.9942 with a standard error of
0.0033, and a perfect solver would add 0.0008 overall).
\texttt{companyTradesAtStockExchange} and \texttt{hasArea} are within noise of
the best posted figures. \texttt{awardWonBy} cannot be evaluated at ten rows.

That leaves \texttt{hasCapacity}, which holds 20.6 percent of the rows and where
the entire field sits below 0.33. We tested our conclusion rather than resting on
it. Forced choice duels are the one instrument available to us that is not a
transformation of the sampled frequency table. We generated every pairwise
comparison among the top six tolerance separated candidates, in both presentation
orders, with eight demonstrations balanced so that the correct answer is at
position A four times and is the larger value four times. Aggregating by Borda
count, with the rule fixed and committed to version control before the pools were
generated, the probe scores 0.3196 on validation against the shipped selector's
0.3608. It loses.

The diagnostics matter more than the verdict, because they distinguish a broken
instrument from a working instrument reporting an absence. The probe is not
degenerate. Position bias is corrected to a mean $P(A)$ of 0.5188, a third of
pairs are decided at a margin above 0.25, and on the 580 validation pairs where
exactly one side is acceptable the model chooses correctly 62.4 percent of the
time, at $z$ of 6.0. The problem is what that accuracy is made of. On the same
pairs, always choosing the larger of the two numbers scores 60.3 percent, and the
probability the model assigns to the larger option rises monotonically with the
ratio between them, from 0.567 below $10^{0.15}$ to 0.727 above $10^{0.75}$. That
is the signature of a magnitude prior, not of recall. The comparative knowledge
above that prior is 2.1 points at $z$ of 0.7.

We then ran the other non-frequency framing, since a duel and a verification have
opposite failure modes: a duel must split its probability mass and so inherits
any position or magnitude prior, while a verifier can accept everything. The
verification probe presents one claim at a time, with balanced yes and no
demonstrations whose false claims are drawn from that demonstration subject's own
pool, and scores a candidate by the log odds of yes against no at the next token.
It is not a yes-bias readout: it accepts only 28.9 percent of candidates at a
positive margin, with a median within row margin spread of 0.250. It scores
0.2680 against 0.3608, a difference of 0.0928 whose 90 percent interval excludes
zero from below, and it is worse inside both strata of a split-half stability
gate. It discriminates, and it discriminates in the wrong direction.

\begin{table}[t]
\centering\small
\begin{tabular}{lr}
\toprule
instrument & val., vs 0.3608 \\
\midrule
cross-frame agreement & 0.3308 vs 0.3308, 0.0000 \\
forced choice duel & 0.3196, $-0.0412$ \\
yes or no verification & 0.2680, $-0.0928$ \\
\bottomrule
\end{tabular}
\caption{Three instruments, each asking the model a different question about the
same candidate set, all agree.}
\label{tab:instruments}
\end{table}

A fourth line of evidence closes it. Sampled frequency is a Monte Carlo estimate
of one functional, so the value of estimating it exactly is the infinite draw
limit of the accuracy curve. Fitting accuracy as $a - b/n$ across pool depths
from 5 to 100 gives a limit of 0.368 to 0.383, a headroom of 0.008 to 0.022 on
the relation. Meanwhile two disjoint fifty draw halves of the same pool disagree
on 38.1 percent of validation rows at almost no cost in accuracy, because the
values they disagree between are equally likely to be wrong.

So the lever is not a better selector, and it is not a better aggregator. On this
relation the model is guessing within a plausible range, and the range happens to
contain the answer often enough to make coverage look like an opportunity.

\section*{Limitations}

Most of the numbers in this paper are validation figures with a standard error of
about 0.021 on the overall metric and 0.03 to 0.05 per relation. Differences
smaller than that, including our apparent margins on \texttt{hasCapacity} and
\texttt{personHasCityOfDeath}, are point estimates and not demonstrated leads. We
say so rather than rounding them into claims.

Thresholds are selected using validation, so validation is a selection set and
not a clean holdout, and the reported validation figure is optimistic by an
amount we estimate at 0.015 to 0.019.

The award relation is not meaningfully evaluable at ten rows and we claim no
result on it. Our validation to test transfer estimates rest on a single edition
and one published system. The coverage analysis measures whether an acceptable
string appears anywhere in a pool, which upper bounds what a selector could
achieve; it is a diagnostic, not a method. And the per relation choice of numeric
aggregator is supported by a confidence interval excluding zero only for
\texttt{hasArea}; for \texttt{hasCapacity} the interval includes zero and the
choice ships as provisional.

\section*{Reproducibility statement}

Code, prompts, seeds, threshold grids, and the parameter audit are available at
\url{https://github.com/thylinao1/lm-kbc-2026}. Pools are regenerable from the
committed configuration, and the manifest hash for every pool consumed by a
submitted predictions file is recorded alongside that file. The full experimental
record, including the runs that failed, is in \texttt{NOTES.local.md} in that
repository.

\bibliography{custom}

\appendix

\section{Negative results and implementation faults in full}
\label{app:negative}

\subsection{Three selection methods that did not work}

\paragraph{Cross-frame consensus loses to the best single frame.} The hypothesis
was that a value surviving a change of register is better evidence than one
merely modal within a single register. On identical cached draws, consensus
scored 0.2887 against 0.3196. The premise was wrong in a specific way: the frames
share one base model and one prior toward round numbers, so their errors are
correlated. For a venue whose true capacity is 5000, three frames independently
proposed 10000, 10500 and 15000. Agreement among correlated predictors certifies
nothing.

\paragraph{Routing by confidence also loses.} Concentration of draws is strongly
predictive of correctness: accuracy rises from 0.000 at support below 0.2 to
0.886 above 0.8. But routing each subject to its most concentrated frame scored
0.2784, and per frame calibrated concentration scored 0.3093, both below the best
single frame. The reason is structural: only about a quarter of frame-subject
pairs reach support above 0.6, and below that every frame sits near 0.15, so the
choice is made precisely where no signal exists. An oracle picking the right
frame per subject would score 0.3814. The information is there and three methods
failed to extract it; we record that as an open problem rather than a result.

\paragraph{Overshoot scaling does not transfer.} A published system reports
multiplying capacity predictions by 0.95. On this substrate all five numeric
frames selected a scale of exactly 1.0 on both train and validation.

\subsection{Three implementation faults}

All three were found by reading errors, not by adding compute, and all three are
the kind of fault that leaves every self-check green.

\paragraph{Integer rounding annihilated small areas.} The numeric writer ended by
rounding to an integer, correct for stadium capacities and destructive for areas:
golds of 0.43, 0.176, 0.063 and 0.008482 square kilometres all became ``0'',
which can never match. Five validation subjects were lost this way, every one
with the correct value already in its pool. Fixing the formatter moved
\texttt{hasArea} from 0.7600 to 0.8500, a gain of 0.019 on the overall metric,
larger than any modelling change we made.

\paragraph{Single-linkage clustering chains.} The published aggregation recipe
clusters draws in log space and takes the median of the largest cluster. With
enough draws, intermediate values bridge distant ones: 10000, 10500, 12000 and
15000 merge into one cluster spanning fifty percent, whose median is a value no
draw proposed and no gold at either end accepts. We replaced it with a fixed
radius count: score each distinct drawn value by how many draws fall within the
grader's own tolerance of it, and take the argmax. A fixed radius cannot chain,
and it optimises exactly the quantity the metric rewards.

\paragraph{A leakage guard correct for one channel and wrong for two.} Our
threshold tuning excludes a channel's own few shot demonstrations from its train
curve, since a demonstration subject has its gold answer sitting in its own
prompt. That guard is sufficient for any single channel measurement and silently
insufficient for any measurement reading two channels at once. Channels draw
different demonstration sets from the same 100 train rows, so a subject clean for
one frame is frequently a demonstration for another, and a cross frame feature
computed on it encodes gold rather than agreement. The union of demonstration
subjects is 64 of 100 train rows for \texttt{hasCapacity} and 74 of 100 for
\texttt{hasArea}, reaching all 100 once the 100 shot frame is included, which
leaves that relation with no leak free train row at all.

We found this by building a cross frame reranker that keeps the shipped frame as
the only proposer of candidates and uses the other frames purely as evidence to
reorder them. Pooled over train and validation it read $+0.0788$ on
\texttt{hasCapacity} and $+0.0476$ on \texttt{hasArea}, with a shuffled label
control far below and every leave one frame out variant positive. Split by split
it was $+0.2206$ on train and exactly $+0.0000$ on validation. On leak free rows
only, 133 for \texttt{hasCapacity} and 126 for \texttt{hasArea}, the effect is
0.0000 with 8 rows up and 8 down, and 2 up and 2 down. Subject grouped cross
validation does not catch this and neither does a label permutation control,
because both test whether the model is learning something real while the leak is
in the feature for that subject. The only check that caught it was asking which
split the gain came from.

\end{document}
